\cotefile{83}
\scanpage{162}
% L1–6
La \deleted{catégorie} donnée d'une var. quadratique
$X$ de $S$, revient à la donnée d'un
\deleted{\unsure{tel syst.|fibré inv.}} $L$ affine et d'un div. quadr. \unsure{dans|\deleted{dans}} $\deleted{\unsure{V}}\check{\mathbb V}(L)$
(lequel sera $X$), ou encore à $L$ et
la fonction de degré $\leq 2$ sur $\check{\mathbb V}(L)$,
dont la partie homogène soit $u\mapsto u^{\otimes2}$.

% L7–13
\deleted{(Dualement, cela \unsure{revient|revi}}
L'algèbre de $X$ \deleted{\unsure{sera|dans}} le dual \inserted{$H\simeq E\otimes L^{-1}$} de $E$,
l'élém$^{\mathrm t}$ unité étant le dual de l'augmentation $E\to k$. \unsure{L$^{k}$a|La} \unsure{transformée|transformation}
par $\otimes L^{-1}$ de $q$ sera une forme
quadratique sur $H$, (\unsure{en|ce} sera la forme
déterminant) égale à $1$ sur $\mathbf{1}$. Il
% L14–21
y aura sur $E$ une unique structure
\deleted{multiplicative} d'alg. quadratique, pour
laquelle $1$ est unité et $q_H$ est le
déterminant. On récupère $\mathrm{Tr}$
\deleted{comme} via la forme bil. associée à $q_H$
comme $x\mapsto\varphi_H(1,x)$, et dual$^{\mathrm{mt}}$ sur
$L$, on récupère $\mathbf{t}\in E$, \deleted{comme} par
la formule
\[
\varphi_{q_E}\bigl(\underset{\substack{\rotatebox{-90}{$\in$}\\L}}{\unsure{u}},\underset{\substack{\rotatebox{-90}{$\in$}\\E}}{x}\bigr)
=\deleted{\unsure{x}}\underset{\substack{\rotatebox{-90}{$\in$}\\L}}{u}\otimes
\underbrace{(\deleted{\unsure{u}}x\wedge\mathbf{t})}_{\in\,\det E\simeq L}\in\deleted{\det E}=L^{\otimes 2}
\]
% L21 lower left
\margin{Attention aux conv. de signe!}
% archive sheet number: 161
